% ==================================================================
% id:       seq-014
% topic:    סדרות
% subtopic: גבולות במובן הרחב
% solution: full
% ==================================================================
\textbf{ עבור כל אחת מהסדרות הבאות, בדקו האם הגבול קיים במובן הרחב ומצאו אותו, או הסבירו מדוע לא קיים:}

\begin{enumerate}\bMs{2}
\item  $\ds{\lim_{n \to \infty} \frac{2 + n^4 \cdot \cos(\pi n)}{2 + n^2}}$
\hint{כתבו את $\cos(\pi n)$ בצורה פשוטה יותר, ובדקו בנפרד את תת-הסדרה של המקומות הזוגיים ואת זו של האי-זוגיים.}
\begin{solbox}
 נסמן את הסדרה ב-
$a_n$.
נזכר כי
$\cos(\pi n)=(-1)^{n}$
לכל
$n$,
ולכן
נציג
$$.a_n=\case{cc}{\frac{2 + n^4}{2 + n^2} & \text{\R{ זוגי }}n\\\frac{2 - n^4}{2 + n^2} & \text{\R{ אי-זוגי }}n}$$
מכיוון ש-
$\lim_{n\to\infty}a_{2n}=\infty$
ו-
$\lim_{n\to \infty}a_{2n+1}=-\infty$,
נסיק כי לא קיים הגבול
במובן הרחב.

\end{solbox}
\item $\ds\lim_{n \to \infty} \frac{\lr{1 + (-1)^n}(n+1)}{3 + 2n}$
\hint{בדקו בנפרד את תת-הסדרה של המקומות הזוגיים ואת זו של האי-זוגיים.}
\begin{solbox}
 נסמן את הסדרה ב-
$a_n$.
נציג
$$.a_n=\case{cc}{\frac{2n + 2}{3 + 2n} &\text{\R{ זוגי }}n\\0 & \text{\R{ אי-זוגי }}n}$$
מכיוון ש-
$\displaystyle\lim_{n\to\infty}a_{2n}=1$
ו-
$\displaystyle\lim_{n\to \infty}a_{2n+1}=0$,
נסיק כי לא קיים הגבול גם במובן הרחב.

הערה: בחישוב הגבול השתמשנו באריתמטיקה של גבולות: הגבול של סכום, של מכפלה ושל מנה של סדרות מתכנסות שווה לסכום, למכפלה ולמנה של הגבולות, כל עוד הגבול של המכנה שונה מ-$0$.
\end{solbox}
\item $\ds\lim_{n \to \infty} \frac{2 + (-1)^n n^2 + 4n^3}{2 + 5n^3 + n^2}$
\hint{חלקו את המונה והמכנה בחזקה הגבוהה ביותר של $n$, ובדקו לאן שואף האיבר שבו מופיע $(-1)^n$.}
\begin{solbox}
     מתקיים
$$.\frac{2 + (-1)^n n^2 + 4n^3}{2 + 5n^3 + n^2}=\frac{\frac{2}{n^3} + \frac{(-1)^n}{n} + 4}{\frac{2}{n^3} + 5 + \frac{1}{n}}$$
מאריתמטיקה של גבולות נסיק כי
$\ds\lim_{n \to \infty} \frac{2 + (-1)^n n^2 + 4n^3}{2 + 5n^3 + n^2}=\frac{4}{5}$.

\end{solbox}
\eMs
\end{enumerate}
